\documentclass{csbulletin}
\selectlanguage{french}
\usepackage[utf8]{inputenc}
\usepackage[all]{nowidow}
\usepackage{comment}
\usepackage{csquotes}
\usepackage{array}
\usepackage{enumitem}
\usepackage{longtable}
% Pied de la 1re page : la classe y ecrit {\doiprefix \doi}, sans argument.
% \doi ne doit donc PAS prendre d'argument ; \doiprefix retire le "DOI:".
\renewcommand{\doi}{Uriel AKAM}
\renewcommand{\doiprefix}{\footnotesize}
\usepackage{makecell}
\usepackage[table]{xcolor}
% Couleur de secours (au cas ou svgnames ne serait pas charge)
\providecolor{Gainsboro}{RGB}{220,220,220}
\usepackage{listings}
\usepackage{textcomp}
\usepackage{upquote}   % apostrophes droites verticales dans les listings
\usepackage{tcolorbox}
\usepackage{fontawesome}
\usepackage{graphicx}  % figures / captures d'ecran
\usepackage{float}     % placement [H]
\usepackage[
  implicit=false,
  hidelinks,
]{hyperref}

% Texte justifie avec de longs mots \texttt insecables : autorise un peu plus
% d'espace entre les mots plutot que de deborder dans la marge.
\setlength{\emergencystretch}{2em}

% Colonne a largeur fixe avec retour a la ligne automatique
\newcolumntype{L}[1]{>{\raggedright\arraybackslash}p{#1}}

% Capture d'ecran : l'image si le fichier existe, sinon un cadre qui signale
% le fichier manquant (le document compile quand meme).
%   \capture{image1.png}  ou  \capture[0.8\linewidth]{image1.png}
\newcommand{\capture}[2][\linewidth]{%
  \IfFileExists{#2}%
    {\includegraphics[width=#1]{#2}}%
    {\fbox{\parbox[c][3cm][c]{\dimexpr#1-2\fboxsep-2\fboxrule\relax}{%
      \centering\small Capture manquante : \texttt{\detokenize{#2}}}}}}

% Configuration des listings (blocs de code)
\lstset{
  basicstyle=\ttfamily\small,
  breaklines=true,
  backgroundcolor=\color{gray!10},
  frame=single,
  rulecolor=\color{gray!40},
  commentstyle=\color{gray},
  showstringspaces=false,
  keepspaces=true,
  upquote=true,
  columns=flexible,
}

% Configuration tcolorbox
\tcbuselibrary{skins}

\begin{document}

\title{Sauvegarde à chaud de la VM \texttt{dev\_mariadb-server2} Rapport de réalisation}
\EnglishTitle{Overlay-based hot backup of the libvirt VM "dev\_mariadb-server2"}
\author{Uriel AKAM}
\podpis{Rapport de réalisation -- Administration Système}
\maketitle
\renewcommand{\figurename}{Figure}
\renewcommand{\tablename}{Tableau}

\begin{abstract}
Ce rapport documente la sauvegarde à chaud de la machine virtuelle
\nolinkurl{dev_mariadb-server2} (hyperviseur libvirt), constituée de deux
disques (\texttt{vda} système, \texttt{vdb} données MariaDB). La méthode repose
sur un \textbf{overlay qcow2 temporaire} par disque : \texttt{virsh
snapshot-create-as --disk-only --quiesce} fige l'écriture sur le disque de base
le temps de la copie, puis \texttt{virsh blockcommit --pivot} réintègre les
écritures accumulées et fait revenir la VM sur son fichier d'origine, sans
interruption de service. Les deux disques ont été traités \textbf{séparément}
pour permettre une copie indépendante vers \nolinkurl{/mnt/backup/libvirt/node3/vm/}.
\end{abstract}


%──────────────────────────────────────────────────────────────────────────
\section{Contexte et objectif}
%──────────────────────────────────────────────────────────────────────────

La VM \texttt{dev\_mariadb-server2} héberge une base MariaDB active et ne peut
être arrêtée pour la sauvegarde. L'objectif est de produire une copie
exploitable des deux disques qcow2 vers l'espace de sauvegarde
\nolinkurl{/mnt/backup/libvirt/node3/vm/}, \emph{sans interrompre} la VM et
\emph{sans copier un fichier en cours d'écriture}.

La contrainte de cohérence impose de ne pas copier directement les fichiers
qcow2 actifs : on bascule temporairement les écritures de la VM vers un
\textbf{overlay}, ce qui fige le fichier de base à un instant précis, le temps
de la copie.


%──────────────────────────────────────────────────────────────────────────
\section{Architecture de la VM}
%──────────────────────────────────────────────────────────────────────────

\begin{longtable}{| L{1.6cm} | L{5.7cm} | L{3.5cm} |}
  \hline
  \rowcolor{Gainsboro!60}
  \makecell{\textbf{Cible}} & \makecell{\textbf{Source}} &
  \makecell{\textbf{Rôle}} \\
  \hline
  \endhead
  \texttt{vda} & \texttt{/var/lib/libvirt/images/} \texttt{dev\_mariadb-server2.qcow2} &
    Disque système \\
  \hline
  \texttt{vdb} & \texttt{/data/vm1/} \texttt{dev\_mariadb-data2.qcow2} &
    Disque de données MariaDB \\
  \hline
\end{longtable}

Destination de la copie : \texttt{/mnt/backup/libvirt/node3/vm/}.


%──────────────────────────────────────────────────────────────────────────
\section{Principe : overlay qcow2 par disque}
%──────────────────────────────────────────────────────────────────────────

Pour chaque disque, on crée un \textbf{overlay} (fichier qcow2 vide adossé au
disque de base via \texttt{backing file}) vers lequel basculent toutes les
nouvelles écritures de la VM. Le fichier de base devient alors \textbf{figé} et
peut être copié en toute sécurité. Une fois la copie terminée, un
\texttt{blockcommit --pivot} réintègre les écritures de l'overlay dans le
fichier de base et fait \textbf{repivoter} la VM dessus. Séquence par disque :
\texttt{snapshot-create-as} $\rightarrow$ copie $\rightarrow$
\texttt{blockcommit --pivot} $\rightarrow$ suppression de l'overlay.

L'option \texttt{--quiesce} fige le système de fichiers invité via
\texttt{qemu-guest-agent}, garantissant une cohérence applicative (important
pour MariaDB).


%──────────────────────────────────────────────────────────────────────────
\section{Procédure exécutée}
%──────────────────────────────────────────────────────────────────────────

Vérification préalable de l'agent invité et de l'état des disques
(lecture seule) :
\begin{lstlisting}[language=bash]
virsh --readonly  list --all  | grep mariadb
\end{lstlisting}

\begin{figure}[H]
  \centering
  \capture{image1.png}
  \caption{Capture d'écran -- vérification préalable  (\texttt{grep}).}
  \label{fig:verif-prealable1}
\end{figure}

\subsection*{Disque \texttt{vda}}


\begin{lstlisting}[language=bash]
virsh --readonly domblklist dev_mariadb-server2
\end{lstlisting}

\begin{figure}[H]
  \centering
  \capture{image3.png}
  \caption{Capture d'écran -- vérification préalable  (\texttt{domblklist}).}
  \label{fig:verif-prealable}
\end{figure}

\subsection*{Disque \texttt{vda}}

\begin{lstlisting}[language=bash]
# 1. Overlay sur vda uniquement (vdb explicitement exclu)
virsh snapshot-create-as --domain dev_mariadb-server2 backup-tmp-vda \
  --diskspec vda,file=/var/lib/libvirt/images/dev_mariadb-server2.backup-overlay.qcow2 \
  --diskspec vdb,snapshot=no \
  --disk-only --atomic --quiesce --no-metadata

# 2. Copie du disque figé
cp --sparse=always /var/lib/libvirt/images/dev_mariadb-server2.qcow2 \
  /mnt/backup/libvirt/node3/vm/dev_mariadb-server2.qcow2

# 3. Recommit + pivot
virsh blockcommit dev_mariadb-server2 vda --active --verbose --pivot

# 4. Nettoyage
rm -f /var/lib/libvirt/images/dev_mariadb-server2.backup-overlay.qcow2
\end{lstlisting}

\begin{figure}[H]
  \centering
  \capture{image2.png}
  \caption{Capture d'écran -- le chemin vers le backup\texttt{vda,vdb}.}
  \label{fig:vda-overlay}
\end{figure}

\subsection*{Disque \texttt{vdb}}

\begin{lstlisting}[language=bash]
# 1. Overlay sur vdb uniquement (vda explicitement exclu)
virsh snapshot-create-as --domain dev_mariadb-server2 backup-tmp-vdb \
  --diskspec vdb,file=/data/vm1/dev_mariadb-data2.backup-overlay.qcow2 \
  --diskspec vda,snapshot=no \
  --disk-only --atomic --quiesce --no-metadata

# 2. Copie du disque figé
cp --sparse=always /data/vm1/dev_mariadb-data2.qcow2 \
  /mnt/backup/libvirt/node3/vm/dev_mariadb-data2.qcow2

# 3. Recommit + pivot
virsh blockcommit dev_mariadb-server2 vdb --active --verbose --pivot

# 4. Nettoyage
rm -f /data/vm1/dev_mariadb-data2.backup-overlay.qcow2
\end{lstlisting}

\begin{figure}[H]
  \centering
  \capture{image3.png}
  \caption{Capture d'écran -- du chemin des disques \texttt{vdb,vda}.}
  \label{fig:vdb-overlay}
\end{figure}

\begin{figure}[H]
  \centering
  \capture{image 4.png}
  \caption{Capture d'écran  et
    \texttt{vdb} : réintégration des écritures de l'overlay et repivotage.}
  \label{fig:blockcommit-pivot}
\end{figure}

\begin{tcolorbox}[colback=yellow!10, colframe=orange!80!black,
                  title={\faExclamationTriangle\ Point de vigilance}]
Les deux disques étant traités \textbf{séparément}, chacun avec son propre
\texttt{--quiesce}, ils sont figés à deux instants distincts. La cohérence
transactionnelle \emph{entre} \texttt{vda} et \texttt{vdb} n'est donc pas
garantie. Pour une cohérence stricte entre les deux disques, un unique
\texttt{snapshot-create-as} regroupant les deux \texttt{--diskspec} reste la
méthode recommandée.
\end{tcolorbox}


%──────────────────────────────────────────────────────────────────────────
\section{Résultats et vérification}
%──────────────────────────────────────────────────────────────────────────

Après \texttt{blockcommit --pivot} sur les deux disques, la VM repointe bien
sur ses fichiers d'origine, sans \texttt{backing file} résiduel :

\begin{lstlisting}[language=bash]
virsh --readonly domblklist dev_mariadb-server2
# Target   Source
# -------------------------------------------------------------
# vda      /var/lib/libvirt/images/dev_mariadb-server2.qcow2
# vdb      /data/vm1/dev_mariadb-data2.qcow2

qemu-img info /var/lib/libvirt/images/dev_mariadb-server2.qcow2   # pas de backing file
qemu-img info /data/vm1/dev_mariadb-data2.qcow2                    # pas de backing file
\end{lstlisting}

\begin{figure}[H]
  \centering
  \capture{image5.png}
  \caption{Capture d'écran -- pour le second}
  \label{fig:verif-finale}
\end{figure}

\begin{tcolorbox}[colback=blue!5, colframe=blue!50!black,
                  title={\faInfoCircle\ Bilan}]
La chaîne \textbf{overlay $\rightarrow$ copie $\rightarrow$ blockcommit/pivot
$\rightarrow$ nettoyage}, exécutée disque par disque, a permis une sauvegarde à
chaud sans interruption de service. La piste d'amélioration principale est le
regroupement des deux disques dans un seul cycle (ou le passage à
\texttt{virsh backup-begin}) afin de garantir une cohérence stricte entre
\texttt{vda} et \texttt{vdb}.
\end{tcolorbox}

\end{document}